\chapter*{Command Quick Reference}
\addcontentsline{toc}{chapter}{Command Quick Reference}\markboth{Command Quick Reference}{}
This appendix consolidates the commands that appear throughout the supplied labs. It is intended as a lookup aid, not a substitute for understanding which VRP view a command belongs to.

\section*{Core VRP and interface commands}
\begin{tabularx}{\linewidth}{>{\ttfamily}p{0.42\linewidth}X}
\toprule
Command & Purpose in the labs\\\midrule
system-view & Enter system configuration view.\\
sysname R1 & Rename the device.\\
interface GigabitEthernet 0/0/0 & Enter an Ethernet interface view.\\
ip address 10.0.13.1 24 & Assign an IPv4 address/prefix.\\
description <text> & Label an interface.\\
display this & Display the configuration in the current view.\\
display interface <interface> & Display detailed interface state.\\
display ip interface brief & Compact IP/interface state summary.\\
ping <address> & Test IP reachability.\\
save & Save the current configuration.\\
display saved-configuration & Display the stored configuration.\\\bottomrule
\end{tabularx}

\section*{Static routing and OSPF}
\begin{tabularx}{\linewidth}{>{\ttfamily}p{0.48\linewidth}X}
\toprule
Command & Purpose\\\midrule
ip route-static 10.0.3.0 24 10.0.23.3 & Add a next-hop static route.\\
ip route-static ... preference 80 & Add a less-preferred backup static route.\\
display ip routing-table & Inspect direct, static, and dynamic routes.\\
ospf 1 router-id 10.0.1.1 & Create/select OSPF process 1 with a router ID.\\
area 0 & Enter OSPF area 0.\\
network 10.0.12.0 0.0.0.255 & Match/advertise the /24 network in the lab.\\
display ospf peer & Detailed OSPF neighbor state.\\
display ospf peer brief & Condensed OSPF neighbor summary.\\\bottomrule
\end{tabularx}

\section*{FTP}
\begin{tabularx}{\linewidth}{>{\ttfamily}p{0.46\linewidth}X}
\toprule
Command & Purpose\\\midrule
ftp server enable & Enable the FTP server.\\
set default ftp-directory flash:/ & Set default server directory.\\
local-user huawei ... & Create/configure the local FTP/PPP user under AAA.\\
display ftp-server & Verify server status and port.\\
ftp 10.0.12.1 & Open an FTP client connection.\\
dir & List files in the FTP session.\\
binary & Set binary transfer mode.\\
get <remote> <local> & Download from server to client.\\
put <local> <remote> & Upload from client to server.\\
bye & Close the FTP session.\\\bottomrule
\end{tabularx}

\section*{DHCP}
\begin{tabularx}{\linewidth}{>{\ttfamily}p{0.48\linewidth}X}
\toprule
Command & Purpose\\\midrule
dhcp enable & Enable DHCP functionality.\\
ip pool pool1 & Create a named global pool.\\
network 10.0.12.0 mask 24 & Define pool network.\\
gateway-list 10.0.12.1 & Define default gateway offered by pool.\\
lease day 1 hour 12 & Configure lease duration.\\
dhcp select global & Use a global pool on an interface.\\
dhcp select interface & Use an interface-based pool.\\
dhcp server excluded-ip-address <IP> & Prevent an address from dynamic allocation.\\
dhcp server dns-list <IP> & Offer a DNS server value.\\
ip address dhcp-alloc & Make a switch Vlanif request an address dynamically.\\
display ip pool name <name> & Inspect pool state and lease counters.\\\bottomrule
\end{tabularx}

\section*{Eth-Trunk, LACP, and VLAN}
\begin{tabularx}{\linewidth}{>{\ttfamily}p{0.51\linewidth}X}
\toprule
Command & Purpose\\\midrule
undo negotiation auto & Disable Ethernet auto-negotiation.\\
speed 100 & Force the lab member-link speed.\\
interface Eth-Trunk 1 & Enter/create the logical bundle.\\
eth-trunk 1 & Add a physical interface to Eth-Trunk 1.\\
mode lacp & Enable static LACP on the Eth-Trunk.\\
lacp priority 100 & Set LACP system/interface priority depending on view.\\
port link-type access & Configure a host-facing access port.\\
port default vlan 4 & Bind an access port to VLAN 4.\\
port link-type trunk & Configure trunk behavior.\\
port trunk allow-pass vlan all & Permit VLANs on the trunk.\\
port link-type hybrid & Configure a hybrid port.\\
port hybrid untagged vlan 2 4 & Send listed VLANs untagged on egress.\\
port hybrid pvid vlan 4 & Classify untagged ingress using VLAN 4 as PVID.\\\bottomrule
\end{tabularx}

\section*{WAN, PPP, PAP, and CHAP}
\begin{tabularx}{\linewidth}{>{\ttfamily}p{0.51\linewidth}X}
\toprule
Command & Purpose\\\midrule
interface Serial 1/0/0 & Enter a serial interface.\\
link-protocol hdlc & Use HDLC encapsulation.\\
baudrate 128000 & Change DCE clock rate in the lab.\\
link-protocol ppp & Use PPP encapsulation.\\
ppp authentication-mode pap & Require PAP authentication.\\
ppp pap local-user huawei password cipher ... & Configure PAP credentials on the authenticated peer.\\
ppp authentication-mode chap & Require CHAP authentication.\\
ppp chap user huawei & Configure CHAP username.\\
ppp chap password cipher ... & Configure CHAP password.\\\bottomrule
\end{tabularx}

\begin{pitfall}
Many VRP commands are view-dependent. If a command is rejected, first check the prompt and confirm that you are in the correct system, interface, protocol, AAA, or FTP view before changing the syntax.
\end{pitfall}
